\documentclass[11pt]{article}
\usepackage[margin=1in]{geometry}
\usepackage{amsmath,amssymb,booktabs,graphicx,float,xurl,array}
\usepackage[hypertexnames=false]{hyperref}
\setlength{\emergencystretch}{6em}
\title{Defect-Reduction Regime Map (Cross-Report, Fixed Full-FPT Fairness)}
\author{Automated benchmark in valley-k-small}
\date{March 9, 2026}
\begin{document}
\sloppy
\maketitle
\section{Scope and Fairness}
Main winner metric compares only full-FPT competitors: sparse exact recursion vs defect-reduced route.
Linear systems are recorded as MFPT reference only.
All claims use fixed-horizon fairness and ratio metric $R=\text{sparse}/\text{defect}$.
Cross-family claims do not mix absolute seconds.
\section{Methods}
Timing protocol: warm-up 1 run, then 3 measured runs with median.
Defect-reduction mode rule: full when defect pairs $\le 120$, otherwise sampled-$z$ extrapolated estimate.
Estimator form: $t_{\mathrm{est}}=t_{\mathrm{setup}}+(t_{\mathrm{eval}}/n_z)\,m+t_{\mathrm{fft}}$.
\section{Main Findings}
Pooled median $R$: 0.0005 (mean 0.0198); probability that the defect-reduced route is faster $P(R>1)$: 0.000.
Global regime decision: No region where the defect-reduced route is faster under fixed-T full-FPT fairness.
Estimation anchors: n=8, median relative error=0.126 (PASS against 25\% threshold).
\section{Figures}
\begin{figure}[H]\centering
\includegraphics[width=0.9\linewidth]{figures/regime_winner_heatmap_two_target.pdf}
\caption{Winner heatmap on two-target workloads.}
\end{figure}
\begin{figure}[H]\centering
\includegraphics[width=0.9\linewidth]{figures/regime_winner_heatmap_reflecting.pdf}
\caption{Winner heatmap on reflecting workloads.}
\end{figure}
\begin{figure}[H]\centering
\includegraphics[width=0.9\linewidth]{figures/regime_speedup_scatter_all.pdf}
\caption{Speedup ratio scatter with full/estimate mode markers.}
\end{figure}
\begin{figure}[H]\centering
\includegraphics[width=0.8\linewidth]{figures/regime_speedup_box_by_T.pdf}
\caption{Distribution of $R$ by fixed horizon $T$.}
\end{figure}
\begin{figure}[H]\centering
\includegraphics[width=0.8\linewidth]{figures/regime_estimation_error_anchor.pdf}
\caption{Full vs estimate relative error on anchor workloads.}
\end{figure}
\begin{figure}[H]\centering
\includegraphics[width=0.98\linewidth]{figures/regime_config_examples.pdf}
\caption{Representative low/medium/high-defect configuration examples.}
\end{figure}
\section{Tables}
\begin{table}[H]
\centering
\small
\caption{Runtime ratio summary by defect bin and horizon.}
\resizebox{\linewidth}{!}{\input{tables/regime_summary_by_bin.tex}}
\end{table}
\begin{table}[H]
\centering
\small
\caption{Anchor workload runtime baselines.}
\resizebox{\linewidth}{!}{\input{tables/regime_anchor_baselines.tex}}
\end{table}
\section{Practical Recommendation Matrix}
\begin{table}[H]
\centering
\small
\begin{tabular}{@{}l>{\raggedright\arraybackslash}p{0.54\linewidth}@{}}
\toprule
Task Type & Recommended Method \\
\midrule
Full FPT under fixed finite horizon & Sparse exact recursion as default baseline; defect-reduced route only when defect support is very small and validated on anchors. \\
MFPT/splitting only & Linear systems (do not use as full-FPT winner metric). \\
Transform-domain analysis with sparse defects & Defect-reduced AW, but keep estimation/full calibration checks. \\
High-defect corridor-like regimes & Prefer sparse recursion; the defect-reduced route often loses after setup/evaluation overhead is counted. \\
\bottomrule
\end{tabular}
\end{table}
\section{Reproducibility}
Core data files: \path{data/manifest.csv}, \path{data/runtime_raw.csv}, \path{data/runtime_summary.json}.
\end{document}
